% Kapittel 1 - Introduksjon
\chapter{Introduksjon}
\label{introduksjon} % Bruk \ref{introduksjon}, for å referere til kapittelet andre steder i oppgaven. 
%----------------------------------------------------------------------------------------

\section{Oppgavetekst}
\label{oppgavetekst}
%Denne oppgaven går ut på å lage en HTML-basert brush editor. Det skal lages en single page web app som kan brukes til å redigere brushfiler. En brushfil er en fil som beskriver hvordan et lakkeringsmønster skal se ut. Filen har mellom 3 og 5 kanaler som tilsier for eksempel hvor mye lakk som skal brukes, hvor mye luft som skal brukes, osv. Web-appen som skal utvikles skal kunne støtte flere ulike typer formfaktorer (PC, tablet, mobil), og må snakke med roboten gjennom et REST API. 

Denne oppgaven går ut på å lage en enkeltside webapplikasjon som kan brukes til å redigere «brush»-filer. En brushfil er en fil som beskriver hvordan et lakkeringsmønster for ABBs malingsroboter skal se ut. Filen bestemmer parametere som hvor mye lakk som skal brukes og sprøytebildet til lakken. Applikasjonen som utvikles skal kunne støtte flere ulike typer formfaktorer (PC, nettbrett og mobiltelefon), og må snakke med roboten gjennom et REST API.


\section{Bakgrunn}
\begin{itemize}
    \item Sammenhengen utstyret/programmet/systemet skal benyttes – motivasjon
    \item Alt som studenten ikke har gjort selv
    \item Teori
\end{itemize}
Kildetest: \cite{Reference1}